\section{View-based Setting}\label{sec:view}
%%%%%%%%%%%%%%%%%%%%%%%%%%%%%%%%%%%%%%%%%%%%%%%%%%%%%%%%%% 

We develop a three-stage model of interaction among a platform, human content creators, and consumers: a content platform hosts a 
\emph{human creator} and an \emph{AI content generator}, both of which compete 
for the attention of a unit mass of \emph{consumers}:
\begin{enumerate}
    \item In Stage~1, the platform chooses a per-view compensation 
rate $r \in \mathbb{R}_+$ paid to the human creator and an algorithmic weight $\beta_H$ 
that determines how prominently human content is promoted relative to 
AI content. 
\item Taking into account these choices, in Stage~2 the human 
creator chooses an effort level that determines the quality of her 
content, while the AI generates content whose quality depends on its 
own baseline capability and what it can learn from the human's 
output. 
\item Finally, in Stage~3, consumers choose between the two providers 
according to a \emph{Hotelling} model. 
\end{enumerate}

Throughout this 
section, we focus on the  \emph{view-based} environment, in 
which both the platform's revenue and the creator's compensation 
scale with realized consumer demand. We solve the game by backward 
induction, beginning with the creator's effort choice and proceeding 
to the platform's design problem.

\subsection{Model}

\subsubsection*{Content Generation and Quality.}

The human creator exerts effort $q_H \geq 0$ (where the subscript $H$ indicates the effort of the human creator), which we identify with 
the quality of her content. Effort is costly, with a 
quadratic cost
\[
C(q_H) \;=\; \tfrac{1}{2}q_H^{2},
\]
which is a standard convex function that captures diminishing returns. The quality of the AI content 
generator depends on both its baseline quality and the human content quality, and is given by 
\begin{equation}\label{eq:qA}
q_A \;=\; \alpha\, q_H + Q,
\end{equation}
where $\alpha \in (0,1)$ governs the AI's  {learning efficiency},
the rate at which it is able to reproduce the human creator's 
quality through training on her content, and $Q \geq 0$ denotes the 
AI's  {baseline quality}, i.e., the quality it can produce 
independently of new human input. Equation ~\eqref{eq:qA} captures 
two features of generative AI in content markets: 
its output improves as humans produce higher-quality 
material that enters its training process, and even without any new human 
input, modern generative systems already produce content of 
non-trivial quality.

\subsubsection*{Consumer Demand.}

We model consumer choice as a one-dimensional Hotelling framework (see, e.g., \cite[Chapter 7]{tirole1988theory}). 
Consumers are distributed uniformly on the unit interval $[0,1]$, 
where the location $x$ indexes a consumer's preference: $x=0$ 
corresponds to a consumer most aligned with the AI, while $x=1$ 
corresponds to a consumer most aligned with the human creator. A 
consumer at location $x$ obtains a utility
\begin{equation}\label{eq:utility}
u_{c}^A(x) \;=\; q_A - t\,x\,\mA, 
\qquad 
u_c^H(x) \;=\; q_H - t(1-x)\,\mH,
\end{equation}
from consuming AI and human content, respectively, where $t > 0$ is 
the typical transportation cost parameter, capturing the unit {mismatch cost} that scales the disutility from 
consuming content misaligned with the consumer's preference. The terms $\mA$ and 
$\mH$ are endogenous mismatch multipliers determined by the 
platform's recommendation algorithm. Letting $\beta_H \in [0,1]$ 
denote the platform's promotion weight on human content and 
$\delta \in (0,1)$ the strength of algorithmic influence on consumer 
search frictions, we have 
\begin{equation}\label{eq:mismatch}
\mA \;=\; 1 - \delta(1-\beta_H), 
\qquad 
\mH \;=\; 1 - \beta_H\delta.
\end{equation}
Both multipliers lie in $[1-\delta,\,1]$. 
Equation~\eqref{eq:mismatch} formalizes the platform's 
role: by choosing a higher $\beta_H$, the platform promotes the human-generated content more compared to the AI-generated content as it increases 
$\mA$ and decreases $\mH$. The parameter $\delta$ caps the 
maximum extent to which algorithmic intervention can impact horizontal differentiation: $\delta \approx 0$ corresponds to 
a neutral platform with no editorial influence, while $\delta \approx 1$ 
corresponds to the maximal algorithmic effect. The symmetric benchmark 
$\beta_H = \tfrac{1}{2}$ yields equal mismatch multipliers and 
serves as a natural reference point.

A consumer at location $x$ chooses AI if $u_c^A(x) \geq 0$ and $u_c^A(x) \geq u_c^H(x)$. Conversely, a consumer at location $x$ chooses the human content if $u_c^H(x) \geq 0$ and $u_c^H(x) \geq u_c^A(x)$. Therefore, the demand for the AI and the human content are 
\begin{equation*}
D_A \triangleq \sup\left\{x \in [0,1]~:~ u_c^A(x) \geq \max\{0, u_c^H(x)\}\right\}, 
~ 
D_H \triangleq \sup\left\{1-x \in [0,1]~:~ u_c^H(x) \geq \max\{0, u_c^A(x)\}\right\}.
\end{equation*}

%Our analysis covers two cases. The first is the uncovered case, $D_A+D_H<1$, where the two providers together do not cover the whole market. Here, the providers do not compete: each acts as a local monopolist on its side of the line and serves only those consumers whose consumption utility is positive. The second is the boundary case, $D_A+D_H=1$, where the market is fully covered. As we show in Section~\ref{ss:vb-equilibrium}, the platform chooses its instruments so that  we always have
% \begin{equation}\label{eq:demand}
% D_A=\frac{q_A}{t\,\mA} \text{ and } D_H=\frac{q_H}{t\,\mH}  \text{ with } \frac{q_A}{t\,\mA}+ \frac{q_H}{t\,\mH} \leq 1.
%  \end{equation}

% $D_A+D_H\leq 1$ at equilibrium; when this constraint binds, the platform deliberately holds total demand at one rather than letting it rise above.

% the human 
% creator if $u_H \geq 0$, and leaves otherwise. We focus on 
% settings in which the joint coverage of the two providers 
% does not exhaust the unit market, so that
% \begin{equation}\label{eq:uncovered}
% \frac{q_A}{t\,\mA} \;+\; \frac{q_H}{t\,\mH} \;\leq\; 1.
% \end{equation}
% Whenever \eqref{eq:uncovered} holds, the providers do not compete with each other: each acts as a local monopolist on its side of the 
% line and serves only those consumers for whom the consumption utility is positive. The resulting 
% demands are the classical Hotelling expressions

% We will see in Section~\ref{ss:vb-equilibrium} that the platform  endogenously chooses its instruments so that \eqref{eq:uncovered}  holds at equilibrium.

\subsubsection*{Creator Utility.}

The platform compensates the human creator on a per-view basis: she 
earns a compensation rate $r \geq 0$ per consumer who chooses her 
content. Given the cost of effort, her utility becomes
\begin{equation}\label{eq:uhV}
u_H^{V}(q_H;\,r,\beta_H) 
\;=\; r\,D_H \;-\; \tfrac{1}{2}q_H^{2}.
\end{equation}

\subsubsection*{Platform Utility.}
We assume that the platform monetizes consumer attention on a 
per-view basis, earning one unit of revenue from each consumer 
served, regardless of whether the consumer chooses AI or human 
content. The pltform pays the creator $r$ for every human view 
but retains the full per-view revenue from AI content. Its 
utility is
\begin{equation}\label{eq:upV}
u_P^{V} \;=\; D_A \;+\; D_H - r D_H 
\;-\; \tfrac{k}{2}\!\left(\beta_H - \tfrac{1}{2}\right)^{2}\!,
\end{equation}
where the quadratic term penalizes deviations of the algorithmic 
weight $\beta_H$ from the neutral benchmark $\tfrac{1}{2}$, with 
$k > 0$ parameterizing the strength of this penalty. The penalty 
captures, in reduced form, the operational, regulatory, and 
reputational costs of moving the recommendation system 
away from neutrality. Three features of \eqref{eq:upV} are worth 
emphasizing. First, the platform values AI views and human views 
asymmetrically: each AI view contributes one unit of  revenue, 
whereas each human view contributes only $1-r$. The platform 
therefore has a preference for shifting consumption 
toward AI content, all else equal. Second, this preference is 
tempered by the recognition that human content sustains the AI's 
own quality through learning, so suppressing human effort harms the platform itself. Third, the algorithmic weight $\beta_H$ enters 
both the demand expressions $D_A$ and $D_H$ through the mismatch 
multipliers in \eqref{eq:mismatch} and the algorithmic-neutrality 
penalty, giving rise to a non-trivial trade-off in the platform's 
optimization problem.

Figure \ref{fig:model:1} illustrates the timing of the game. 

%%%
\begin{figure}
    \centering
\begin{tikzpicture}[
    >=Stealth,
    stage_label/.style={font=\footnotesize\sffamily, fill=white, inner sep=2pt},
    player_box/.style={rounded corners=4pt, draw=#1, thick, fill=#1!6,
                        text width=3.8cm, align=left, inner sep=8pt, font=\small},
    timeline_dot/.style={circle, fill=black!70, inner sep=2.2pt},
    action_arrow/.style={->, thick, #1, shorten >=2pt, shorten <=2pt},
]
 
% ── Timeline ──────────────────────────────────────────────
\draw[->, line width=1.2pt, black!60]
    (-0.5, 0) -- (14.5, 0);
 
% Stage markers
\foreach \x/\lab in {1/Stage 1, 6.5/Stage 2, 12/Stage 3} {
    \node[timeline_dot] at (\x, 0) {};
    \node[stage_label, below=4pt] at (\x, 0) {\lab};
}
 
% ── Platform (Stage 1, top-left) ──────────────────────────
\node[player_box=violet] (platform) at (1, 1.8) {%
    \textbf{\textsf{Platform chooses}}\\[2pt]
    Compensation rate:\; $r$\\
    Promotion weight:\; $\beta_H$
};
 
\draw[action_arrow=violet] (platform.south) -- (1, 0.15);
 
% ── Human Creator (Stage 2, bottom-center) ────────────────
\node[player_box=red!70!black] (creator) at (6.5, -1.8) {%
    \textbf{\textsf{Human creator chooses}}\\[4pt]
    Content quality:\; $q_H$
};
 
\draw[action_arrow=red!70!black] (creator.north) -- (6.5, -0.15);
 
% ── Consumers (Stage 3, top-right) ────────────────────────
\node[player_box=blue!70!black] (consumers) at (12, 1.8) {%
    \textbf{\textsf{Consumers choose}}\\[4pt]
    between AI \& Human\\
    generated content
};
 
\draw[action_arrow=blue!70!black] (consumers.south) -- (12, 0.15);
 
% ── Annotations along the arrow ──────────────────────────
\node[font=\scriptsize\itshape, text=black!50, above=6pt] at (3.75, 0)
    {$r,\;\beta_H$ observed};
\node[font=\scriptsize\itshape, text=black!50, above=6pt] at (9.25, 0)
    {$q_H,\;q_A$ observed};
\end{tikzpicture}
\caption{The timing of the game among the platform, the human creator, and the consumers}
    \label{fig:model:1}
\end{figure}
%%%

\subsection{Equilibrium Characterization}\label{ss:vb-equilibrium}

We solve the game by backward induction. 


\subsubsection*{Stage 3: Consumer's Choice and Demand Regimes}

In the final stage, consumers choose between AI-generated and human-generated content according to a Hotelling framework, as described above. The following proposition characterizes the consumer's choice. 
%%%
\begin{proposition}\label{pro:ConsumerChoice:View-Based}
Given transportation cost $t>0$, mismatch multipliers $m_A$, $m_H$, and content qualities $q_A$, $q_H$, the consumer demand for AI content ($D_A$) and human content ($D_H$) is determined by two distinct regimes:
\begin{enumerate}
    \item \textbf{Uncovered Regime:} If $\frac{q_A}{t m_A} + \frac{q_H}{t m_H} < 1$, the market is not fully covered. Each provider acts as a local monopolist, and demands are given by:
    \begin{equation*}
        D_A = \frac{q_A}{t m_A}, \quad D_H = \frac{q_H}{t m_H}
    \end{equation*}
    \item \textbf{Covered Regime:} If $\frac{q_A}{t m_A} + \frac{q_H}{t m_H} \ge 1$, the market is fully covered ($D_A + D_H = 1$). The demands are determined by an indifferent consumer located at $x^*$, yielding:
    \begin{equation*}
        D_A = \frac{q_A - q_H + t m_H}{t(m_A + m_H)}, \quad D_H = \frac{q_H - q_A + t m_A}{t(m_A + m_H)}
    \end{equation*}
\end{enumerate}
\end{proposition}
Proposition \ref{pro:ConsumerChoice:View-Based} highlights a key difference between uncovered and covered markets. In the uncovered regime, each provider effectively serves a separate segment of consumers without directly competing for the same marginal users. As a result, demand for each type of content depends only on its own quality (higher quality or a lower mismatch cost expands the set of consumers willing to consume that content). In particular, demand increases monotonically with own quality and is independent of the competitor's quality.

By contrast, in the covered regime, all consumers choose either AI or human content, so the two providers compete for the same viewers. Demand is therefore determined by the location of the indifferent consumer and depends on relative rather than absolute quality. An increase in the quality of AI content not only raises demand for AI content but also reduces demand for human content, and vice versa. Similarly, lower mismatch costs shift consumers toward the corresponding provider by making that type of content more appealing relative to the alternative.


% Proposition \ref{pro:ConsumerChoice:View-Based} establishes that when the market is uncovered, the demand for both human and AI contents depends only on their quality and mismatch multipliers and, in particular, is increasing in their quality. This changes when the market is covered because there will be competition and therefore the demand for human content creators (and similarly, the AI content generator) depends on the relative quality of human versus AI content. 

%Stage~2 determines the creator's  best-response effort given the platform's instruments; Stage~1 then optimizes those instruments, anticipating the creator's response.


\subsubsection*{Stage 2: Creator's Optimal Effort}

Given the platform's choices of $r$ and $\beta_H$, the human creator selects effort $q_H \ge 0$ to maximize her utility $u_H = r D_H - \frac{1}{2} q_H^2$. Because consumer demand changes when the market becomes fully covered, the creator's marginal benefit of effort depends on the coverage regime, as characterized above, and we obtain the following characterization. 

\begin{proposition}\label{pro:Optimal:q_H:View-Based}
For any compensation rate $r>0$, algorithmic weight $\beta_H \in [0,1]$, mismatch parameter $t>0$, and AI parameters $\alpha \in (0,1)$ and $Q \ge 0$, assume $t(1-\delta)>Q$ such that the AI does not trivially monopolize the market. With corresponding algorithmic multipliers $m_A = 1-\delta(1-\beta_H)$ and $m_H = 1-\delta\beta_H$, let $\overline{q}_H = \frac{m_H(t m_A - Q)}{m_A + \alpha m_H}$ be the boundary effort level where the market exactly covers. 

The human creator's unique best-response effort $q_H^*$ is given by:
\begin{enumerate}
    \item \textbf{Uncovered Optimum:} $q_H^* = \frac{r}{t m_H}$ if $r \le \frac{t m_H^2 (t m_A - Q)}{m_A + \alpha m_H}$.
    \item \textbf{Covered Optimum:} $q_H^* = \frac{r(1-\alpha)}{t(m_A + m_H)}$ if $r \ge \frac{t(m_A + m_H)m_H(t m_A - Q)}{(1-\alpha)(m_A + \alpha m_H)}$.
    \item \textbf{Boundary Optimum:} $q_H^* = \overline{q}_H$ for intermediate values of $r$.
\end{enumerate}
\end{proposition}
%%%
Proposition \ref{pro:Optimal:q_H:View-Based} shows that the human creator's optimal choice of quality depends critically on whether additional quality expands the market or merely reallocates consumers away from the AI competitor. When the compensation rate is low, the creator has weak incentives to invest in quality and chooses an effort level that keeps the market uncovered. In this regime, higher quality directly increases the creator's own share of the market by attracting previously uncovered consumers, so the optimal effort balances this market-expansion benefit against the cost of quality provision.

As the compensation rate increases, the creator is willing to invest more in quality until the market becomes fully covered. Beyond this point, further quality improvements no longer expand total demand because every consumer is already consuming some form of content. Instead, additional quality only shifts viewers from AI content toward human content. Here, the losses from the competition effect balance the gains from the market-expansion effect, and consequently, for an intermediate range of compensation rates, the creator optimally remains exactly at the coverage boundary $\overline{q}_H$. However, when compensation becomes sufficiently large, it becomes worthwhile to compete aggressively for consumers in the covered market, leading the creator to increase quality beyond the boundary level.

% Proposition \ref{pro:Optimal:q_H:View-Based} establishes that for a small compensation rate, the human creator chooses a quality that results in an uncovered market. In contrast, a large enough compensation rate incentivizes the human content creator to improve its quality, resulting in a covered market. For intermediary ranges of the compensation rate, the gain from a higher pay exactly balances the loss from losing consumers to the competitor AI content generator, and therefore, the human content quality remains fixed. 


% the creator's best-response effort is dependent on whether the market is covered.  When the market is uncovered, the creator only competes with the outside option. Every unit of effort brings in new users, so she works hard based on her pay $r$ and her own mismatch cost. However, when the market is covered, the game changes. Now, she must steal users from the AI. Because the AI learns from her effort at rate $\alpha$, working harder also makes the AI better. This learning effect hurts her ability to steal users, so her marginal benefit of effort drops by a factor of $1-\alpha$. This sudden drop in incentives creates the boundary regime. For middle values of $r$, the creator stops exactly where the market covers. Pushing past this point is not worth the cost, because she would face an AI that gains from her improvements.


\subsubsection*{Stage 1: Platform's Optimal Design}

Now we solve the first stage. The platform jointly finds the best pay rate $r$ and the best algorithm weight $\beta_H$ to maximize its profit. Because the mismatch terms $m_A = 1 - \delta(1-\beta_H)$ and $m_H = 1 - \delta\beta_H$ depend on the algorithm, the best pay rate also depends on $\beta_H$. We solve this joint problem by first finding the optimal $r$ for any given $\beta_H$, and then plugging it back into the profit function to find the optimal $\beta_H$.


% {\color{red}PLEASE ADD the assumptions and then the proposition we talked about here (see main4), and then have a discussion of it, and then say in the appendix BLAH, we characterize the platform's optimal decision with a relaxed assumption on $k$ ... }



Sufficiently large values of the penalty parameter $k$ guarantee concavity of the platform's optimization problem. Motivated by this result, we impose the following assumption throughout the analysis.

\begin{assumption}\label{ass:k_large}
The penalty $k$ is sufficiently large such that the platform's utility is concave in $\beta_H$.
\end{assumption}

Under Assumption~\ref{ass:k_large}, the platform admits an interior optimal solution. We therefore focus the main analysis on this case. We characterize the concavity lower bound for $k$ in Lemma \ref{lem:k_bar:View-Based} in the appendix. Also, the complementary setting in which $k$ is small, leading to corner solutions with $\beta_H^* \in {0,1}$, is discussed in the appendix.


\begin{assumption}\label{ass:nonsat}
   Neither side can saturate the market in isolation:
   \[
   Q < t\,\mA \quad \text{and} \quad 1 < t\,\mH 
   \qquad \text{for all } \beta_H \in [0,1].
   \]
   Equivalently, $Q < t(1-\delta)$ and $t > 1/(1-\delta)$.
\end{assumption}

Assumption~\ref{ass:nonsat} ensures that the AI's standalone demand 
$Q/(t\mA)$ is strictly less than 1, so that the existence of a human 
creator is consequential: at $q_H = 0$ the market is genuinely 
under-served, and the platform has a non-trivial reason to incentivize 
human effort. Symmetrically, the condition $t\,\mH > 1$ ensures that 
each unit of creator effort generates less than one unit of demand: 
since $D_H = q_H/(t\,\mH)$, the human's reach grows in effort at the 
rate $1/(t\,\mH) < 1$. The existence of an AI content generator is 
therefore consequential: at any effort level the platform can 
incentivize, the human alone leaves the market under-served, and the 
platform has a non-trivial reason to host AI-generated content 
alongside her. 



\begin{proposition}[Platform's Joint Optimal Design]
\label{prop:optimal:r_beta:View-Based}
Given the framework parameters $\alpha, Q, \delta, k$ and under Assumptions \ref{ass:k_large}-\ref{ass:nonsat}, the platform's joint optimization over compensation $r$ and algorithmic promotion $\beta_H$ reduces to a single-variable maximization over $\beta_H \in [0,1]$. 

\begin{enumerate}
    \item \textbf{Optimal Compensation Rate:} For any given promotion weight $\beta_H$, the optimal compensation rate $r^*(\beta_H)$ is uniquely clamped at either the uncovered optimum or the boundary regime:
    \begin{equation}
        r^*(\beta_H) = \min\left( \frac{m_A + \alpha m_H}{2m_A}, \; \frac{t m_H^2 (t m_A - Q)}{m_A + \alpha m_H} \right)
    \end{equation}
    
    \item \textbf{Optimal Algorithmic Weight:} Let $M(\beta_H)$ represent the derivative (marginal revenue) of the platform's utility without the penalty term. $M(\beta_H)$ evaluates piece-wise conditionally on $\beta_H$:
    \begin{equation}
        M(\beta_H) = 
        \begin{cases}
            \delta \left[ \frac{1}{2t^2} \left( \frac{1}{m_H} + \frac{\alpha}{m_A} \right) \left( \frac{1}{m_H^2} - \frac{\alpha}{m_A^2} \right) - \frac{Q}{t m_A^2} \right] & \text{if } (m_A + \alpha m_H)^2 \le 2 t m_A m_H^2 (t m_A - Q), \\[10pt]
            \frac{2\delta m_H (t m_A - Q)}{(m_A + \alpha m_H)^3} \left[ t m_A^2 - \alpha t m_H^2 - Q(2-\delta) \right] & \text{otherwise.}
        \end{cases}
    \end{equation}
    The optimal algorithmic weight is then formally given by the implicit bounded mapping:
    \begin{equation}
        \beta_H^* = \max\left(0, \min\left(1, \; \frac{1}{2} + \frac{M(\beta_H^*)}{k} \right)\right)
    \end{equation}
\end{enumerate}
\end{proposition}

The platform's choice of $r$ balances two opposing forces. Raising $r$ 
is costly because the platform pays it on every human view, but it 
also induces the creator to exert more effort, which both expands her 
own demand and improves the AI's content through 
learning. The interior optimum $r^{\text{unc}}$ exceeds 
$\tfrac{1}{2}$ precisely because of this AI-learning spillover: human 
effort is worth more to the platform than the creator's own revenue 
contribution would suggest, since better human content also makes the 
AI's content better at no extra cost. The boundary rate 
$r^{\text{bdy}}$ becomes binding when the market is small 
relative to the providers' joint reach, when mismatch costs are low, 
the AI's baseline quality is high, or learning is fast. In that case, 
paying the creator more would only push the market past saturation, 
redistributing already-captured consumers between the two providers 
while raising the platform's payment to the creator, so the platform deliberately 
under-incentivizes the creator to keep the market exactly at the 
coverage boundary.



\begin{corollary}
\label{cor:beta_0.5:View-Based}
    The optimal algorithmic weight $\beta_H^*$ is larger than $1/2$ if and only if 
    \[
    Q \;<\; \frac{1-\alpha}{2}\cdot\min\!\left\{\frac{1+\alpha}{tm_0},\; tm_0\right\}
    \]
where $tm_0 = 1-\delta/2$.
\end{corollary}

A human-favoring promotion is justified by {small} baseline AI
quality, with the threshold tightening as AI learning improves.

{\color{red} Will add the intuition for the choice of $\beta_H$}

